\section{Construction and elementary structure of \(X\)}

In this section we construct the graph

\[
X=\KC(L(P))
\]

in several equivalent ways and record its elementary graph-theoretic
properties. These descriptions will be used interchangeably throughout the
remainder of the paper.

\subsection{Canonical-cover and tensor-product models}

Let

\[
Y=L(P)
\]

be the line graph of the Petersen graph.

The principal object of this paper is the canonical bipartite double cover

\[
X=\KC(Y).
\]

Equivalently,

\[
X\cong Y\times K_2,
\]

where \(K_2\) denotes the complete graph on two vertices and the product is
the tensor product of graphs.

Thus

\[
V(X)
=
V(Y)\times\Ftwo,
\]

and two vertices

\[
(u,\varepsilon),
\qquad
(v,\delta)
\]

are adjacent if and only if

\[
u\sim_Y v
\]

and

\[
\delta
=
\varepsilon+1
\pmod2.
\]

The projection

\[
\pi:
X
\longrightarrow
Y,
\qquad
(u,\varepsilon)
\longmapsto
u,
\]

is a regular two-fold covering projection whose deck involution is

\[
\tau(u,\varepsilon)
=
(u,\varepsilon+1).
\]

\begin{proposition}

The graph \(X\) is connected, bipartite, quartic, and satisfies

\[
|V(X)|=30,
\qquad
|E(X)|=60.
\]

\end{proposition}

\begin{proof}

Since

\[
|V(Y)|=15,
\]

the double cover contains

\[
2\cdot15=30
\]

vertices.

Every vertex of \(Y\) has degree four, and the covering projection is
locally bijective. Hence every vertex of \(X\) also has degree four.

Applying the Handshake Lemma,

\[
|E(X)|
=
\frac{30\cdot4}{2}
=
60.
\]

By Proposition~2.2 the canonical bipartite double cover of a connected
non-bipartite graph is connected and bipartite. Since \(L(P)\) contains
triangles, the result follows.

\end{proof}

The remaining subsections describe the same graph from alternative
viewpoints that will be useful in later sections.


\subsection{The Petersen-edge model}

A vertex of the line graph \(Y=L(P)\) corresponds to an edge of the
Petersen graph. Consequently every vertex of

\[
X=\KC(L(P))
\]

may be written in the form

\[
(e,\varepsilon),
\]

where

\[
e\in E(P),
\qquad
\varepsilon\in\Ftwo.
\]

Thus \(X\) consists of two copies of the fifteen edges of the Petersen
graph.

Two vertices

\[
(e,\varepsilon),
\qquad
(f,\delta)
\]

are adjacent precisely when

\begin{enumerate}
\item the edges \(e\) and \(f\) of the Petersen graph share a common
vertex, and

\item

\[
\delta=\varepsilon+1
\pmod2.
\]

\end{enumerate}

Equivalently,

\[
(e,\varepsilon)
\sim
(f,\varepsilon+1)
\]

if and only if

\[
e\cap f\neq\varnothing.
\]

This coordinate system will be used throughout the paper. In particular,
the explicit vertex labels appearing in the computational certificates are
obtained by fixing a labeling of the fifteen edges of the Petersen graph,
as recorded in Appendix~A.


\subsection{The neighborhood-incidence model}

Let

\[
\mathcal N(Y)
\]

denote the neighborhood incidence structure of the graph \(Y\). Its points
are the vertices of \(Y\), and its blocks are the open neighborhoods

\[
N_Y(v),
\qquad
v\in V(Y).
\]

Since \(Y\) is simple and quartic, every block has cardinality four.

The Levi graph of this incidence structure has one bipartition consisting
of the vertices of \(Y\) and the other consisting of the neighborhood
blocks. A vertex is incident with a block precisely when it belongs to the
corresponding neighborhood.

\begin{proposition}

The graph \(X=\KC(L(P))\) is isomorphic to the Levi graph of the
neighborhood incidence structure of \(Y=L(P)\).

\end{proposition}

\begin{proof}

The canonical bipartite double cover separates each vertex of \(Y\) into
two sheets while preserving every adjacency relation between neighboring
vertices.

The neighborhood-incidence construction records exactly the same incidence
information by replacing one sheet with neighborhood blocks. Identifying
each neighborhood block with its corresponding lifted fiber produces a
graph isomorphism.

\end{proof}

The canonical-cover, tensor-product, Petersen-edge, and neighborhood-
incidence descriptions therefore define the same graph. We shall pass
freely between these models according to convenience.


\subsection{Connectedness, bipartiteness, and girth}

By Proposition~2.2, the graph

\[
X=\KC(L(P))
\]

is connected and bipartite.

The remaining elementary invariant is its girth.

\begin{proposition}

The girth of \(X\) is six.

\end{proposition}

\begin{proof}

Since \(X\) is bipartite, every cycle has even length.
It therefore suffices to exclude cycles of length four.

Suppose

\[
x_1x_2x_3x_4x_1
\]

were a 4-cycle in \(X\).

Projecting to the base graph \(Y=L(P)\) yields vertices

\[
u_1,u_2,u_3,u_4
\]

such that

\[
u_i\sim u_{i+1}
\]

for each \(i\), with indices taken modulo four.

Because the covering projection is locally bijective, the projected walk is
a closed walk of length four in \(Y\).

However, \(L(P)\) contains no 4-cycles.

Indeed, every pair of adjacent vertices in \(L(P)\) corresponds to two
incident edges of the Petersen graph, and the Petersen graph contains no
triangles or squares from which a 4-cycle in its line graph could arise.

Hence no 4-cycle exists in \(X\).

Since \(L(P)\) contains triangles, each triangle lifts to a cycle of length
six in the canonical bipartite double cover.

Therefore the shortest cycle in \(X\) has length six.

\end{proof}

Consequently,

\[
X
\]

is a connected bipartite quartic graph of girth six.


\subsection{Symmetry and connectivity}

The canonical bipartite double cover inherits every automorphism of the
base graph.

Every automorphism

\[
g\in\Aut(L(P))
\]

acts on

\[
X=\KC(L(P))
\]

by

\[
(u,\varepsilon)
\longmapsto
(g(u),\varepsilon).
\]

This lifted action preserves adjacency and commutes with the deck
involution

\[
\tau(u,\varepsilon)
=
(u,\varepsilon+1).
\]

Consequently,

\[
\Aut(L(P))
\times
\langle\tau\rangle
\cong
S_{5}\times C_{2}
\]

acts faithfully on \(X\).

The full automorphism group of \(X\) is determined later in
Section~6. At present we record only the elementary consequences of the
lifted action.

\begin{proposition}

The graph \(X\) is vertex-transitive, edge-transitive, and arc-transitive.

\end{proposition}

\begin{proof}

The line graph of the Petersen graph is arc-transitive.

Since every automorphism lifts to the canonical bipartite double cover,
the lifted action remains transitive on vertices, edges, and oriented
edges.

\end{proof}


\subsection{Metric and spectral profile}

The graph \(X\) has diameter four.

Its distance distribution from any vertex is

\[
(1,4,12,11,2),
\]

showing that every vertex has

\[
1+4+12+11+2
=
30
\]

vertices at distances

\[
0,1,2,3,4,
\]

respectively.

The adjacency matrix of the canonical bipartite double cover has block
form

\[
A_X
=
\begin{pmatrix}
0&A_Y\\
A_Y&0
\end{pmatrix},
\]

where \(A_Y\) is the adjacency matrix of the line graph \(Y=L(P)\).

Consequently every eigenvalue

\[
\lambda
\]

of \(A_Y\) gives rise to the pair of eigenvalues

\[
\pm\lambda
\]

of \(A_X\).

The resulting spectrum is

\[
\operatorname{Spec}(X)
=
\{
4^{1},
2^{10},
1^{4},
(-1)^{4},
(-2)^{10},
(-4)^{1}
\}.
\]

These metric and spectral invariants will be used later only as recognition
data; the proofs of the principal theorems do not depend upon them.


\subsection{Recognition profile}

The elementary invariants established in this section are summarized in
Table~\ref{tab:recognition-profile}.

\begin{table}[h]
\centering
\begin{tabular}{ll}
\toprule
Invariant & Value\\
\midrule
Order & \(30\)\\
Size & \(60\)\\
Degree & \(4\)\\
Connected & Yes\\
Bipartite & Yes\\
Girth & \(6\)\\
Diameter & \(4\)\\
Distance distribution & \((1,4,12,11,2)\)\\
Spectrum &
\(\{4^{1},2^{10},1^{4},(-1)^{4},(-2)^{10},(-4)^{1}\}\)\\
\bottomrule
\end{tabular}
\caption{Elementary graph-theoretic invariants of
\(X=\KC(L(P))\).}
\label{tab:recognition-profile}
\end{table}

These invariants provide an elementary recognition profile for the graph.
The subsequent sections investigate the additional structure arising from
its quotient frames, homogeneous-space realization, and automorphism
group.
